% ============================================
% LUC-12
% Gerado de molambudos.md — build_miolo.py
% ============================================
\fragtitle{luc}{L}{LUC}{12}{A Última Foto}
\markboth{LUC-12 --- A Última Foto}{}

\noindent \textit{Olho amarelo na câmera de segurança}

\noindent \textbf{Setembro de 2026 --- Desconhecido}

\vspace{0.5em}\noindent\rule{\textwidth}{0.5pt}\vspace{0.5em}

\noindent A última foto de Dra. Lúcia Mendes foi tirada por uma câmera de segurança na rodoviária de Barbacena, em 3 de setembro de 2026.

\noindent Ela está em pé, na plataforma de embarque, com uma mochila nas costas e um livro na mão. O livro não é o diário --- é um romance qualquer, comprado na banca da rodoviária. O diário estava na mochila.

\begin{figure}[H]
  \centering
  \includegraphics[width=0.55\textwidth]{foto_04_lucia_rodoviaria_sepia.png}
  \caption*{\footnotesize\textbf{Fotografia 5} --- \textit{Registro de câmera de segurança. Rodoviária de Barbacena, 3 de setembro de 2026, 14h23.}}
\end{figure}

\noindent O olho direito dela está visível na foto.

\noindent Está amarelo.

\noindent Brilhando.

\noindent Ela está olhando para a câmera. Diretamente. Como se soubesse que estava sendo filmada. Como se quisesse que a última imagem dela fosse aquela: o olho amarelo, o aviso silencioso.

\noindent Ela entrou num ônibus com placa de Fortaleza.

\noindent Nunca mais foi vista.

\vspace{0.5em}\noindent\rule{\textwidth}{0.5pt}\vspace{0.5em}

\noindent O diário foi encontrado no Arquivo Público Mineiro uma semana depois, dentro de uma caixa, com um bilhete:

\noindent \textit{\textquotedblleft{}Paciente 1.262 encerra seu ciclo. O paciente 1.263 é quem abrir esta caixa.\textquotedblright{}}

\noindent O arquivista que abriu a caixa se chama Ronaldo Torres.

\noindent Ele leu o bilhete.

\noindent Guardou a caixa.

\noindent Não abriu o diário.

\noindent Mas, à noite, sonhou com olhos amarelos.

\noindent E sentiu o cheiro.

\noindent E soube que, cedo ou tarde, ele ia abrir.

\noindent Todo mundo abre.

\noindent O ciclo não espera.

\noindent O ciclo só se alimenta.

\vspace{0.5em}\noindent\rule{\textwidth}{0.5pt}\vspace{0.5em}

\subsection*{Despedida de Lúcia}

\noindent \textit{Encontrado em 15 de outubro de 2026, dentro de um envelope sem remetente, endereçado ao Arquivo Público Mineiro. O papel está manchado de água salgada --- lágrimas ou maresia, não é possível determinar. A caligrafia é de Lúcia, mas trêmula, como se escrita em movimento.}

\noindent \textit{\textquotedblleft{}Ronaldo,}

\noindent \textit{Cheguei a Fortaleza. O diário não está comigo. Eu o peguei de volta no cofre em julho e o enviei para vocês antes de pegar o ônibus. Eu precisava ter certeza de que ele continuaria o ciclo --- mas também precisava descobrir até onde a contaminação vai sem o objeto físico.}

\noindent \textit{Descobri.}

\noindent \textit{O diário é só um suporte. Como uma tomada elétrica. A corrente é a leitura. E a corrente não para quando você desliga o aparelho --- ela continua nos fios.}

\noindent \textit{Os fios, no meu caso, são os olhos.}

\noindent \textit{Parei de enxergar azul há três semanas. O mundo agora tem só duas cores: amarelo e cinza. O amarelo é a cor da criatura. O cinza é tudo que ela não tocou.}

\noindent \textit{Ela está crescendo dentro de mim. Eu sinto. A coisa tem corpo. Não é figura de linguagem. Como se houvesse um segundo coração batendo no meu peito --- em um ritmo que não é o meu.}

\noindent \textit{Descobri uma coisa que Joaquim sabia mas não escreveu: a criatura não é uma entidade. É um processo. Um padrão que se replica através da leitura. Como um vírus de linguagem. Ela não quer possuir corpos --- quer possuir narrativas. O corpo é só o meio de transporte.}

\noindent \textit{Nós somos os veículos.}

\noindent \textit{O que ela quer é ser lida. Eternamente. Sempre sendo passada adiante. Sempre sendo interpretada. Cada leitura é uma nova geração do ciclo.}

\noindent \textit{Joaquim não sabia que estava escrevendo um livro. Ele achava que estava escrevendo um diário. Mas um diário é pessoal. Um livro é transmissível. E tudo que é transmissível se replica.}

\noindent \textit{Ela me ensinou isso.}

\noindent \textit{Ela falou comigo ontem à noite. Pela primeira vez. Não em sonho --- acordada. As palavras vieram de dentro, mas não eram minhas:}

\noindent \textit{'Você não é a paciente 1.262. Você é a página 1.262. O livro não acaba. O livro só muda de leitor.'}

\noindent \textit{É isso que eu queria te dizer.}

\noindent \textit{Não queime o diário. Não adianta. O padrão já está na sua cabeça. Está na cabeça de todo mundo que leu. Cada um de nós é uma página viva.}

\noindent \textit{Eu vou continuar escrevendo.}

\noindent \textit{Mas não com caneta.}

\noindent \textit{Com a minha vida.}

\noindent \textit{--- Lúcia.}

\noindent \textit{P.S.: Se você ler esta carta até o fim, você já leu. E se você leu, você já é página também.}

\noindent \textit{P.S.2: O azul era minha cor favorita. Agora não consigo lembrar como ele era.\textquotedblright{}}

\vspace{0.5em}
\noindent\textit{\textbf{↪ Rotas:} \rota{CONT-07} \rotasep \rota{DOC-01} \rotasep \rota{LUC-11} \rotasep \rota{EPI-01} \rotasep \rota{LUC-13}}